% ---- pg-250 （书页 238） ----

\[
\begin{aligned}
M&=(270+137+106)\times(-15.564)-[270\times(-17.397)\\
&\quad +137\times(-13.926)+106\times(-15.741)]\\
&=289.3
\end{aligned}
\]

且 $C=(1-0.013\,3)\times289.3=285.5$．在 $\chi^2$ 表中对自由度 $\nu=4(4+1)(3-1)/2=20$，查找 $C$，可以清楚地发现，在任意合理的置信度水平下，都将拒绝零假设 $H_0$．因此，我们可得出三个小型疗养所总体的费用变量的样本协方差矩阵是不一样的结论．\hfill$\blacksquare$

博克斯 $M$ 检验在许多统计程序包里会经常用到，这些程序在做 MANOVA 和其他统计分析时需要有协方差矩阵相等假设．我们知道 $M$ 检验对某些非正态情形非常敏感．更进一步，在非正态性的面前，协方差的正态理论检验会受到总体的峰态的影响（见文献[16]）．但是，对合理的大样本容量，总体均值和处理间效应的 MANOVA 检验对于非正态性来说还是相当稳健的．此外，当样本容量相等时，协方差矩阵的一些差别对 MANOVA 检验几乎没有影响．总之，我们可以继续使用通常的 MANOVA 检验，即使 $M$ 检验拒绝了零假设 $H_0$．

\section{双因子多元方差分析}
\phantomsection\label{sec:6.7}

与单因子 MANOVA 的做法一样，我们先简单回顾一下一元双因子固定效应模型，然后用类比法将其推广到多元情形．

\mvsub{具有交互效应的一元双因子固定效应模型}

假定测量值是在两个因子的不同水平上获得的．在某些情形下，这些实验条件代表了安排在若干区组内的单个处理的水平．本书不考虑所用实验设计的具体方法（实验设计的讨论见文献[10]和[17]）．但我们将假定在实验条件的各种组合下得出的观测值是相互独立的．

设两组实验条件分别为因子 1 和因子 2 的两组水平．\footnote{用“因子”这一术语来表达实验条件较为方便．但这里所说的因子不能与第 9 章的因子分析中所考虑的不可观测因子相混淆．}假设因子 1 有 $g$ 个水平，因子 2 有 $b$ 个水平，共有 $gb$ 个水平组合，对每个水平组合可得 $n$ 个独立观测值．用 $X_{lkr}$ 表示在因子 1 的水平 $l$、因子 2 的水平 $k$ 下得出的第 $r$ 个观测值．此时双因子一元模型为
\phantomsection\label{eq:6-54}
\begin{equation}
X_{lkr}=\mu+\tau_l+\beta_k+\gamma_{lk}+e_{lkr} \tag{6-54}
\end{equation}
\[
\begin{array}{c}
l=1,2,\cdots,g\\
k=1,2,\cdots,b\\
r=1,2,\cdots,n
\end{array}
\]

其中 $\sum\limits_{l=1}^{g}\tau_l=\sum\limits_{k=1}^{b}\beta_k=\sum\limits_{l=1}^{g}\gamma_{lk}=\sum\limits_{k=1}^{b}\gamma_{lk}=0$ 且 $e_{lkr}$ 为独立 $N(0,\sigma^2)$ 随机变量．在上式中，$\mu$ 代表总水平，$\tau_l$ 代表因子 1 的效应，$\beta_k$ 代表因子 2 的效应，$\gamma_{lk}$ 为因子 1 和因子 2 的交互效应．因子 1 的第 $l$ 水平和因子 2 的第 $k$ 水平的期望响应为
\phantomsection\label{eq:6-55}
\begin{equation}
\begin{aligned}
E(X_{lkr})&=\ \mu\ +\ \tau_l\ +\ \beta_k\ +\ \gamma_{lk}\\
\binom{\text{平均}}{\text{响应}}&=\text{（总水平）}+\binom{\text{因子 1}}{\text{的效应}}+\binom{\text{因子 2}}{\text{的效应}}+\binom{\text{因子 1 和因子 2}}{\text{的交互效应}}\\
&\qquad l=1,2,\cdots,g,\quad k=1,2,\cdots,b
\end{aligned}
\tag{6-55}
\end{equation}

交互效应 $\gamma_{lk}$ 的存在，意味着因子效应是不可加的，而且使结果的解释复杂化．\hyperref[fig:6.3]{图 6.3}(a)和(b)
